\documentclass[12pt,a4paper]{article}
\usepackage[utf8]{inputenc}
\usepackage{amsmath,amsfonts,amssymb}
\usepackage{listings}
\usepackage{xcolor}
\usepackage{geometry}
\usepackage{hyperref}
\usepackage{fancyhdr}

\geometry{margin=1in}
\pagestyle{fancy}
\fancyhf{}
\rhead{Pandas - Complete Guide}
\lfoot{Python for Data Science}
\rfoot{\thepage}

\lstset{
    language=Python,
    basicstyle=\ttfamily\small,
    keywordstyle=\color{blue}\bfseries,
    commentstyle=\color{green!60!black},
    stringstyle=\color{red},
    numbers=left,
    numberstyle=\tiny\color{gray},
    stepnumber=1,
    numbersep=8pt,
    showstringspaces=false,
    breaklines=true,
    frame=single,
    backgroundcolor=\color{gray!10}
}

\title{\textbf{Pandas - Complete Guide}}
\author{Python for Data Science - Beginner's Level}
\date{\today}

\begin{document}

\maketitle
\tableofcontents
\newpage

\section{Overview}

Pandas is a powerful data manipulation and analysis library built on NumPy. It provides DataFrame and Series objects for handling structured data.

\section{Installation and Import}

\begin{lstlisting}
# Installation
# pip install pandas

# Import
import pandas as pd
import numpy as np
\end{lstlisting}

\section{Core Data Structures}

\subsection{Series}

\begin{lstlisting}
# Create Series
s1 = pd.Series([1, 2, 3, 4, 5])
s2 = pd.Series([1, 2, 3], index=['a', 'b', 'c'])
s3 = pd.Series({'A': 1, 'B': 2, 'C': 3})

# Series properties
print(s2.values)      # [1 2 3]
print(s2.index)       # Index(['a', 'b', 'c'])
print(s2.dtype)       # int64
print(s2.shape)       # (3,)
\end{lstlisting}

\subsection{DataFrame}

\begin{lstlisting}
# Create DataFrame
data = {
    'Name': ['Alice', 'Bob', 'Charlie', 'Diana'],
    'Age': [25, 30, 35, 28],
    'City': ['NYC', 'LA', 'Chicago', 'Boston'],
    'Salary': [50000, 60000, 70000, 55000]
}
df = pd.DataFrame(data)

# DataFrame properties
print(df.shape)       # (4, 4)
print(df.columns)     # Index(['Name', 'Age', 'City', 'Salary'])
print(df.index)       # RangeIndex(start=0, stop=4, step=1)
print(df.dtypes)      # Data types of each column
\end{lstlisting}

\section{Data Selection and Indexing}

\subsection{Basic Selection}

\begin{lstlisting}
# Column selection
names = df['Name']              # Series
subset = df[['Name', 'Age']]    # DataFrame

# Row selection
first_row = df.iloc[0]          # First row by position
named_row = df.loc[0]           # First row by label

# Slicing
first_three = df.iloc[0:3]      # First 3 rows
columns_slice = df.iloc[:, 1:3] # All rows, columns 1-2
\end{lstlisting}

\subsection{Advanced Selection}

\begin{lstlisting}
# Boolean indexing
high_salary = df[df['Salary'] > 55000]
young_high_earners = df[(df['Age'] < 30) & (df['Salary'] > 50000)]

# Query method
filtered = df.query('Age > 25 and Salary < 65000')

# isin method
cities_filter = df[df['City'].isin(['NYC', 'LA'])]

# String methods
name_filter = df[df['Name'].str.startswith('A')]
\end{lstlisting}

\section{Data Manipulation}

\subsection{Adding and Modifying Data}

\begin{lstlisting}
# Add new column
df['Bonus'] = df['Salary'] * 0.1
df['Full_Info'] = df['Name'] + ' - ' + df['City']

# Modify existing data
df.loc[df['Age'] > 30, 'Category'] = 'Senior'
df.loc[df['Age'] <= 30, 'Category'] = 'Junior'

# Apply functions
df['Age_Group'] = df['Age'].apply(lambda x: 'Young' if x < 30 else 'Mature')
\end{lstlisting}

\subsection{Removing Data}

\begin{lstlisting}
# Drop columns
df_no_bonus = df.drop('Bonus', axis=1)
df_subset = df.drop(['Bonus', 'Full_Info'], axis=1)

# Drop rows
df_no_first = df.drop(0, axis=0)
df_filtered = df.drop(df[df['Age'] < 25].index)

# Remove duplicates
df_unique = df.drop_duplicates()
df_unique_names = df.drop_duplicates(subset=['Name'])
\end{lstlisting}

\section{Data Cleaning}

\subsection{Handling Missing Data}

\begin{lstlisting}
# Create data with missing values
df_missing = df.copy()
df_missing.loc[1, 'Salary'] = np.nan
df_missing.loc[2, 'Age'] = np.nan

# Check for missing data
print(df_missing.isnull().sum())
print(df_missing.info())

# Handle missing data
df_dropped = df_missing.dropna()                    # Drop rows with any NaN
df_filled = df_missing.fillna(0)                    # Fill with 0
df_mean = df_missing.fillna(df_missing.mean())      # Fill with mean
\end{lstlisting}

\subsection{Data Type Conversion}

\begin{lstlisting}
# Convert data types
df['Age'] = df['Age'].astype('int32')
df['Salary'] = df['Salary'].astype('float64')

# Convert to datetime
dates = pd.Series(['2023-01-01', '2023-02-01', '2023-03-01'])
df['Date'] = pd.to_datetime(dates)

# Convert to categorical
df['City'] = df['City'].astype('category')
\end{lstlisting}

\section{Data Analysis}

\subsection{Descriptive Statistics}

\begin{lstlisting}
# Basic statistics
print(df.describe())                    # Summary statistics
print(df['Salary'].mean())              # Mean salary
print(df['Age'].median())               # Median age
print(df['City'].value_counts())        # Count by city

# Correlation
numeric_df = df.select_dtypes(include=[np.number])
correlation_matrix = numeric_df.corr()
\end{lstlisting}

\subsection{Grouping and Aggregation}

\begin{lstlisting}
# Group by single column
city_groups = df.groupby('City')
print(city_groups['Salary'].mean())     # Average salary by city
print(city_groups.size())               # Count by city

# Group by multiple columns
df['Age_Group'] = pd.cut(df['Age'], bins=[0, 30, 40, 100], 
                        labels=['Young', 'Middle', 'Senior'])
multi_group = df.groupby(['City', 'Age_Group'])
print(multi_group['Salary'].agg(['mean', 'count']))

# Custom aggregations
agg_result = df.groupby('City').agg({
    'Salary': ['mean', 'max', 'min'],
    'Age': 'mean'
})
\end{lstlisting}

\subsection{Pivot Tables}

\begin{lstlisting}
# Create pivot table
pivot = df.pivot_table(
    values='Salary',
    index='City',
    columns='Age_Group',
    aggfunc='mean',
    fill_value=0
)

# Cross-tabulation
crosstab = pd.crosstab(df['City'], df['Age_Group'])
\end{lstlisting}

\section{Data Merging and Joining}

\subsection{Concatenation}

\begin{lstlisting}
# Create additional data
df2 = pd.DataFrame({
    'Name': ['Eve', 'Frank'],
    'Age': [32, 29],
    'City': ['Seattle', 'Miami'],
    'Salary': [65000, 58000]
})

# Concatenate vertically
combined = pd.concat([df, df2], ignore_index=True)

# Concatenate horizontally
df_extra = pd.DataFrame({'Department': ['IT', 'HR', 'Finance', 'Marketing']})
horizontal = pd.concat([df, df_extra], axis=1)
\end{lstlisting}

\subsection{Merging}

\begin{lstlisting}
# Create related data
departments = pd.DataFrame({
    'Name': ['Alice', 'Bob', 'Charlie', 'Diana'],
    'Department': ['IT', 'HR', 'Finance', 'Marketing']
})

# Inner join
merged_inner = pd.merge(df, departments, on='Name', how='inner')

# Left join
merged_left = pd.merge(df, departments, on='Name', how='left')
\end{lstlisting}

\section{File I/O Operations}

\subsection{Reading Data}

\begin{lstlisting}
# CSV files
df_csv = pd.read_csv('data.csv')
df_csv_custom = pd.read_csv('data.csv', sep=';', header=0, index_col=0)

# Excel files
df_excel = pd.read_excel('data.xlsx', sheet_name='Sheet1')

# JSON files
df_json = pd.read_json('data.json')
\end{lstlisting}

\subsection{Writing Data}

\begin{lstlisting}
# CSV files
df.to_csv('output.csv', index=False)
df.to_csv('output_custom.csv', sep=';', header=True)

# Excel files
df.to_excel('output.xlsx', sheet_name='Data', index=False)

# JSON files
df.to_json('output.json', orient='records')
\end{lstlisting}

\section{Time Series Analysis}

\begin{lstlisting}
# Create time series data
dates = pd.date_range('2023-01-01', periods=100, freq='D')
ts_data = pd.DataFrame({
    'Date': dates,
    'Value': np.random.randn(100).cumsum()
})
ts_data.set_index('Date', inplace=True)

# Time-based selection
jan_data = ts_data['2023-01']
week_data = ts_data['2023-01-01':'2023-01-07']

# Resampling
monthly = ts_data.resample('M').mean()      # Monthly average
weekly = ts_data.resample('W').sum()        # Weekly sum

# Rolling operations
ts_data['Rolling_Mean'] = ts_data['Value'].rolling(window=7).mean()
\end{lstlisting}

\section{Performance Tips}

\begin{lstlisting}
# Use vectorized operations
df['Total'] = df['Quantity'] * df['Price']  # Fast

# Use appropriate data types
df['Category'] = df['Category'].astype('category')  # Memory efficient

# Use query for complex filtering
result = df.query('Price > 100 and Category == "Electronics"')  # Fast

# Use chaining for multiple operations
result = (df
          .query('Age > 25')
          .groupby('City')['Salary']
          .mean()
          .sort_values(ascending=False))
\end{lstlisting}

\section{Summary}

Pandas provides:
\begin{itemize}
    \item \textbf{DataFrame/Series}: Powerful data structures
    \item \textbf{Data cleaning}: Handle missing values and duplicates
    \item \textbf{Data analysis}: Grouping, aggregation, and statistics
    \item \textbf{File I/O}: Read/write various formats
    \item \textbf{Time series}: Date/time operations and resampling
    \item \textbf{Merging}: Join and combine datasets
    \item \textbf{Performance}: Optimized operations for large datasets
\end{itemize}

Essential for data analysis and manipulation in Python!

\end{document}